
For each state $\ket{x_0,x_1}$ for $x_i\in \{0,1\}$, apply the following circuit:

\[
    \begin{quantikz}
    \lstick{} & \gate{H} & \ctrl{1} & \gate{H} & \qw \\
    \lstick{} & \gate{H} & \targ{}  & \gate{H} & \qw
    \end{quantikz}
\]
Repeat the calculation with the following circuit:
\[
    \begin{quantikz}
        \lstick{} & \targ{} & \qw \\
        \lstick{} & \ctrl{-1} & \qw
    \end{quantikz}
\]
Conclude that these two unitary transformations are equal, and that the intuition that the
C-NOT gate does not change the control qubit is false, when we work in the ${\ket{+}, \ket{-}}$ basis.

\subsection*{Solution}
Let us define:
\begin{align*}
    H_2 &= H \otimes H\\
    C &= CNOT\\
    \tilde{C} &= CNOT_{2\rightarrow 1}\\
\end{align*}
Where $H$ is Hadamard operator, $CNOT$ is controlled-not where first qubit is control and second qubit is target. $CNOT_{2\rightarrow 1}$ is controlled-not where second qubit is control and first is the target.
Applying the first circuit is equivalent to performing the $H_2 C H_2$ operator:\\

For state $\ket{00}$:
\begin{align*}
    H_2 C H_2\ket{00}
    &=H_2 C\ket{++}\\
    &=\frac{1}{2}H_2 C(\ket{00}+\ket{01}+\ket{10}+\ket{11})\\
    &=\frac{1}{2}H_2(\ket{00}+\ket{01}+\ket{11}+\ket{10})\\
    &=H_2\ket{++}\\
    &=\ket{00}
\end{align*}

For state $\ket{01}$:
\begin{align*}
    H_2 C H_2\ket{01}
    &=H_2 C\ket{+-}\\
    &=\frac{1}{2}H_2 C(\ket{00}-\ket{01}+\ket{10}-\ket{11})\\
    &=\frac{1}{2}H_2(\ket{00}-\ket{01}+\ket{11}-\ket{10})\\
    &=\frac{1}{2}H_2(\ket{0}(\ket{0}-\ket{1})+\ket{1}(\ket{1}-\ket{0}))\\
    &=\frac{1}{2}H_2(\ket{0}-\ket{1}(\ket{0}-\ket{1})\\
    &=H_2\ket{--}\\
    &=\ket{11}
\end{align*}

For state $\ket{10}$:
\begin{align*}
    H_2 C H_2\ket{10}
    &=H_2 C\ket{-+}\\
    &=\frac{1}{2}H_2 C(\ket{00}+\ket{01}-\ket{10}-\ket{11})\\
    &=\frac{1}{2}H_2(\ket{00}+\ket{01}-\ket{11}-\ket{10})\\
    &=H_2\ket{-+}\\
    &=\ket{10}
\end{align*}

For state $\ket{11}$:
\begin{align*}
    H_2 C H_2\ket{01}
    &=H_2 C\ket{+-}\\
    &=\frac{1}{2}H_2 C(\ket{00}-\ket{01}-\ket{10}+\ket{11})\\
    &=\frac{1}{2}H_2(\ket{00}-\ket{01}-\ket{11}+\ket{10})\\
    &=\frac{1}{2}H_2(\ket{0}(\ket{0}-\ket{1})-\ket{1}(\ket{1}-\ket{0}))\\
    &=\frac{1}{2}H_2(\ket{0}+\ket{1}(\ket{0}-\ket{1})\\
    &=H_2\ket{+-}\\
    &=\ket{01}
\end{align*}

The second circuit is equivalent to performing the $\tilde{C}$ operator, this by definition gives us the following behavior:
\begin{align*}
    \tilde{C}\ket{00} &= \ket{00}\\
    \tilde{C}\ket{01} &= \ket{11}\\
    \tilde{C}\ket{10} &= \ket{10}\\
    \tilde{C}\ket{11} &= \ket{01}\\
\end{align*}
Comparing the results of the two circuits we observe the same behavior for all 4 base states, therefore concluding they are equal.
Furthermore, observing (for example) the behavior of the first circuit on the state $\ket{01}$ we notice that the following is true:
\[
    C\ket{+-} = \ket{--}
\]
This demonstrates a case in which, when operating in the $\{+,-\}$ basis, C-NOT changes the control (first) qubit